% TEMPLATE-MILC-v3.tex - plantilla maestra UNICA de las guias MILC v3.
%
% La usan las tres familias de guias, cada una con su driver:
%   · Grados 6-11  -> scripts/build-guias-g11.py
%   · Bebras       -> scripts/build-guias-bebras.py
%   · Semillero    -> scripts/build-guias-semillero.py
%
% Antes habia tres copias casi identicas (~400 lineas iguales, 6 distintas):
% cada mejora habia que aplicarla tres veces, y en la practica no se hacia
% (el motor de recursos vivia solo en dos de ellas). Lo que varia entre
% familias esta parametrizado en 6 placeholders que cada driver llena
% (se nombran aqui SIN su sintaxis real: el builder reemplaza por texto plano,
% tambien dentro de los comentarios, y un valor multilinea rompe el preambulo):
%
%   HEADER_IZQ        encabezado izquierdo de pagina
%   HEADER_DER        encabezado derecho de pagina
%   PORTADA_ETIQUETA  rotulo sobre el numero grande (GRADO/MOMENTO/MODULO)
%   PORTADA_SERIE     linea de serie de la portada
%   PORTADA_ID        identificador de la guia en la portada
%   PORTADA_CIRCULO   el circulito de grado (°), o vacio si no aplica
%
% El resto de claves las documenta content/guias/_SCHEMA.md. El script valida
% que no queden placeholders sin reemplazar antes de escribir el archivo final.

\documentclass[11pt,letterpaper]{article}
\usepackage[margin=2.0cm]{geometry}
\usepackage[table]{xcolor}
\usepackage{fontspec}
\usepackage[spanish]{babel}
\usepackage{tikz}
% Librerías TikZ para los diagramas de las guías (bloque `recursos:`):
%  · babel  → babel en español vuelve activos caracteres como «>», y TikZ los usa
%             en sus opciones (>=stealth, ->). Sin esta librería, cualquier
%             diagrama con flechas revienta con «\language@active@arg> has an
%             extra }». Debe ir ANTES de que se dibuje nada.
%  · positioning        → sintaxis «below left=2pt and 3pt».
%  · decorations.pathreplacing → llaves ({brace}) para señalar rangos.
\usetikzlibrary{babel,positioning,decorations.pathreplacing}
\usepackage{graphicx}
% Circuitos: el trabajo de electrónica se hace hoy con micro:bit y sus sensores
% integrados (MakeCode, sin cableado), así que no se necesitan esquemáticos.
% Si algún día entran montajes con protoboard, basta descomentar esta línea:
% los diagramas .tex/.tikz declarados en `recursos:` ya se incrustan solos.
% \usepackage{circuitikz}
\usepackage[most]{tcolorbox}
\usepackage{tabularx}
\usepackage{array}
\usepackage{fancyhdr}
\usepackage{titlesec}
\usepackage[document]{ragged2e}  % bandera a la derecha en todo el cuerpo
\usepackage{enumitem}
\usepackage{microtype}

% Helvetica Neue no trae la flecha U+2192, asi que cada «→» escrito literalmente
% en el YAML se compilaba a NADA, en silencio (462 flechas en 61 guias: rutas del
% tipo «paleta → tipografia → lockup» salian rotas). Mapeamos el caracter a su
% version matematica, que si tiene glifo. Asi el autor puede seguir escribiendo
% «→» comodamente en el YAML.
\usepackage{newunicodechar}
\newunicodechar{→}{\ensuremath{\rightarrow}}
% Mismo problema con los simbolos de marca y vinetas: el «✓» de «una palomita ✓»
% salia como cuadrito vacio en el PDF de todas las guias que lo usaban.
\usepackage{pifont}
\newunicodechar{✓}{\ding{51}}
\newunicodechar{✗}{\ding{55}}
\newunicodechar{✔}{\ding{52}}
\newunicodechar{•}{\textbullet}
\newunicodechar{≥}{\ensuremath{\geq}}
\newunicodechar{≤}{\ensuremath{\leq}}

\setmainfont{Helvetica Neue}
\setsansfont{Helvetica Neue}
\setlength{\parindent}{0pt}
\setlength{\parskip}{6pt}
% Sin particion de palabras: en 6.º-7.º una palabra cortada por guion al final
% de linea («Comuni-cacion») frena la lectura mucho mas que una linea corta.
\hyphenpenalty=10000
\exhyphenpenalty=10000
\pretolerance=2000
\tolerance=3000
\setlength{\emergencystretch}{3em}
\linespread{1.10}
\renewcommand{\arraystretch}{1.25}
\setlist{leftmargin=5mm,itemsep=1mm,topsep=1mm}
\newcolumntype{Y}{>{\RaggedRight\arraybackslash}X}

\definecolor{milcMagenta}{HTML}{E6007E}
\definecolor{milcVino}{HTML}{5A0038}
\definecolor{milcTurquesa}{HTML}{0093A5}
\definecolor{milcVerde}{HTML}{58A923}
\definecolor{milcMostaza}{HTML}{E5B400}
\definecolor{milcOcre}{HTML}{B86600}
\definecolor{milcNegro}{HTML}{241816}
\definecolor{milcGris}{HTML}{F7F7F4}
\definecolor{milcLinea}{HTML}{E8E2DD}
\definecolor{milcGradePrimary}{HTML}{0D9488}
\definecolor{milcGradeDark}{HTML}{115E59}
\definecolor{milcGradeSoft}{HTML}{CCFBF1}
\definecolor{milcGradeAccent}{HTML}{CFE7FF}
\definecolor{milcGradeText}{HTML}{FFFFFF}
\color{milcNegro}

\pagestyle{fancy}
\fancyhf{}
\lhead{\small Territorio interior · Educación socioemocional}
\rhead{\small Grado 10° · Momento 8 · MILC v3}
\cfoot{\small\thepage}
\renewcommand{\headrulewidth}{0pt}

\newtcolorbox{titlebox}[1]{
  enhanced,colback=#1,colframe=#1,arc=7mm,boxrule=0pt,
  left=7mm,right=7mm,top=3mm,bottom=3mm,
  before skip=8pt,after skip=8pt,
  fontupper=\sffamily\bfseries\Large\color{white}
}

\newtcolorbox{softbox}[3]{
  enhanced,colback=#2,colframe=#1,borderline west={4pt}{0pt}{#1},
  arc=2mm,boxrule=.4pt,left=4mm,right=4mm,top=3mm,bottom=3mm,
  before skip=6pt,after skip=8pt,title={#3},coltitle=white,
  colbacktitle=#1,fonttitle=\sffamily\bfseries,fontupper=\RaggedRight,
  attach boxed title to top left={xshift=3mm,yshift=-2mm},
  boxed title style={arc=2mm, boxrule=0pt}
}

\newtcolorbox{drawbox}[2]{
  enhanced,colback=white,colframe=#1,arc=4mm,boxrule=1pt,
  left=5mm,right=5mm,top=4mm,bottom=4mm,before skip=8pt,after skip=8pt,
  title={#2},coltitle=white,colbacktitle=#1,fonttitle=\sffamily\bfseries,
  fontupper=\RaggedRight,
  attach boxed title to top left={xshift=4mm,yshift=-2mm},
  boxed title style={arc=3mm, boxrule=0pt}
}

\newtcolorbox{infoband}[1]{
  enhanced,colback=#1!8,colframe=#1!50,arc=2mm,boxrule=.5pt,
  left=4mm,right=4mm,top=2mm,bottom=2mm,before skip=4pt,after skip=4pt,
  fontupper=\small\RaggedRight
}

% Puente narrativo entre secciones (contrato editorial Conéctate).
% Da continuidad: "Acabas de X. Ahora vas a Y porque Z."
% Visualmente: banda izquierda en color del grado, texto en itálica pequeña.
\newtcolorbox{puentebox}{
  enhanced,colback=milcGradeSoft,colframe=milcGradeDark,
  borderline west={3pt}{0pt}{milcGradeDark},
  arc=0mm,boxrule=0pt,left=5mm,right=4mm,top=2.5mm,bottom=2.5mm,
  before skip=4pt,after skip=8pt,
  fontupper=\itshape\small\color{milcNegro}\RaggedRight
}
% La flecha va en modo matematico a proposito: Helvetica Neue NO tiene el glifo
% U+2192, asi que \textrightarrow se compilaba a NADA («Missing character: There
% is no → in font Helvetica Neue») y el puente salia sin flecha. En modo
% matematico la flecha viene de la fuente matematica y si se dibuja.
\newcommand{\puente}[1]{%
  \begin{puentebox}%
    \textcolor{milcGradeDark}{$\rightarrow$}\hspace{2mm}#1%
  \end{puentebox}%
}

\newcommand{\linead}[1]{\noindent\color{milcLinea}\rule{#1}{0.5pt}\color{milcNegro}}
\newcommand{\respuesta}[1]{\par\vspace{2mm}\foreach \n in {1,...,#1}{\noindent\color{milcLinea}\rule{\linewidth}{0.5pt}\par\vspace{5mm}}\color{milcNegro}}
\newcommand{\checkbox}{\textcolor{milcGradeDark}{\fbox{\phantom{x}}}\hspace{5pt}}

% ─── Listas aireadas para las guías socioemocionales ────────────────────
% Componer las verificaciones como «\checkbox texto\\» deja el texto que salta
% de línea debajo del cuadrito y apelmaza el bloque. Estos entornos alinean la
% continuación y airean los ítems. Los usa Territorio Interior; cualquier
% familia puede hacerlo.
\newlist{checklista}{itemize}{1}
\setlist[checklista]{
  label=\textcolor{milcGradeDark}{\fbox{\phantom{x}}},
  leftmargin=9mm, labelsep=3.2mm, itemsep=2.4mm,
  topsep=2.5mm, parsep=0pt, partopsep=0pt, before=\RaggedRight
}
\newlist{pasoslista}{enumerate}{1}
\setlist[pasoslista]{
  label=\textcolor{milcGradeDark}{\sffamily\bfseries\arabic*.},
  leftmargin=8mm, labelsep=3mm, itemsep=2.8mm,
  topsep=1mm, parsep=0pt, partopsep=0pt, before=\RaggedRight
}

\newcommand{\phasecard}[4]{
\begin{tcolorbox}[enhanced,colback=#3,colframe=#2,arc=3mm,boxrule=.5pt,height=3.05cm,valign=center,halign=center,left=1.5mm,right=1.5mm,top=2mm,bottom=2mm]
{\sffamily\bfseries\fontsize{22}{24}\selectfont\textcolor{#2}{#1}}\par
{\sffamily\bfseries\small #4}
\end{tcolorbox}
}

% ── Piezas del rediseno 2026 (legibilidad 6.º-11.º) ───────────────────
% Banda de estacion: reemplaza el titlebox macizo. Titulo a la izquierda,
% dato practico (tiempo/modalidad) a la derecha, en una sola linea.
\newcommand{\milcestacion}[2]{%
\begin{tcolorbox}[enhanced,colback=#1,colframe=#1,arc=2mm,boxrule=0pt,
  left=5mm,right=5mm,top=2.6mm,bottom=2.6mm,before skip=11pt,after skip=7pt]
{\sffamily\bfseries\fontsize{15}{18}\selectfont\color{white}#2}
\end{tcolorbox}}

% Tarjeta de paso numerada: sustituye las tablas «Pilar / Aplicacion», que
% eran formato de consulta docente, no de estudiante. El numero va en un
% circulo sobre el borde para que la secuencia se lea de un vistazo.
\newcommand{\milcpaso}[3]{%
\begin{tcolorbox}[enhanced,colback=white,colframe=milcLinea,arc=1.5mm,boxrule=.9pt,
  left=14mm,right=4mm,top=3mm,bottom=3mm,before skip=4pt,after skip=4pt,
  fontupper=\RaggedRight,
  overlay={\node[anchor=north west,circle,fill=#2,text=white,inner sep=0pt,
    minimum size=7.4mm,font=\sffamily\bfseries\small]
    at ([xshift=3.6mm,yshift=-3.2mm]frame.north west) {#1};}]
#3
\end{tcolorbox}}

% Casilla marcable de tamano util con lapiz (la anterior era \fbox{\phantom{x}},
% de alto variable segun el texto que la seguia).
\newcommand{\milcmarca}{\framebox[3.8mm]{\rule{0pt}{3.4mm}}}

% Fila marcable de lista suelta (checks de la estacion 1).
\newcommand{\milccheckrow}[1]{%
\par\vspace{1.8mm}\noindent
\makebox[6.4mm][l]{\raisebox{-.55mm}{\textcolor{milcNegro}{\milcmarca}}}%
\parbox[t]{\dimexpr\linewidth-6.4mm\relax}{\RaggedRight #1}\par}

% Celda marcable dentro de la rubrica (tabla de autoevaluacion). Misma
% sangria francesa, pero pensada para vivir dentro de una columna X.
\newcommand{\milcopcion}[1]{%
\makebox[5.2mm][l]{\raisebox{-.55mm}{\textcolor{milcNegro}{\milcmarca}}}%
\parbox[t]{\dimexpr\linewidth-5.2mm\relax}{\RaggedRight #1}}

% ── Recursos visuales de la guía ──────────────────────────────────────
% Declarados en el bloque `recursos:` del YAML y emitidos por el builder.
% \guiaFigura   → imagen rasterizada o PDF (foto, carta celeste, montaje).
% \guiaDiagrama → fuente TikZ/circuitikz (.tex) incrustada como vector:
%                 es la vía para circuitos y esquemáticos editables.
% #1 = ruta relativa al .tex · #2 = pie (puede ir vacío)
\newcommand{\guiaFigura}[2]{%
\begin{center}
\includegraphics[width=0.88\linewidth,height=0.38\textheight,keepaspectratio]{#1}\\[1.5mm]
{\footnotesize\itshape\textcolor{milcNegro!75}{#2}}
\end{center}
}

\newcommand{\guiaDiagrama}[2]{%
\begin{center}
\input{#1}\\[1.5mm]
{\footnotesize\itshape\textcolor{milcNegro!75}{#2}}
\end{center}
}

\begin{document}
\thispagestyle{empty}

% ── Portada ───────────────────────────────────────────────────────────
% Antes: un rectangulo de color a pagina completa. Se veia bien en pantalla,
% pero en la fotocopiadora del colegio se comia el toner y salia gris sucio.
% Ahora la banda ocupa el tercio superior y el resto va en blanco.
\begin{tikzpicture}[remember picture,overlay]
  \path[fill=milcGradePrimary]
    (current page.north west) rectangle ([yshift=-10.6cm]current page.north east);

  \node[anchor=north west,text=milcGradeText] at ([xshift=2.0cm,yshift=-1.55cm]current page.north west)
    {\sffamily\bfseries\fontsize{12}{14}\selectfont MOMENTO};
  \node[anchor=north west,text=milcGradeText,opacity=.85] at ([xshift=2.0cm,yshift=-2.35cm]current page.north west)
    {\sffamily\bfseries\fontsize{8.5}{10}\selectfont TERRITORIO INTERIOR · SOCIOEMOCIONAL};
  \node[anchor=north east,text=milcGradeText,opacity=.85,align=right] at ([xshift=-2.0cm,yshift=-1.55cm]current page.north east)
    {\sffamily\bfseries\fontsize{9}{12}\selectfont I.E. Sor María Juliana\\Cartago, Valle del Cauca};

  \node[anchor=west,text=milcGradeText] (gradonum) at ([xshift=1.85cm,yshift=-7.1cm]current page.north west)
    {\sffamily\bfseries\fontsize{112}{112}\selectfont 8};
  % El circulito de grado (°) solo aplica a las guias de grado («11°»); en
  % Bebras y Semillero el numero grande es un momento/modulo, no un grado.
  % Se posiciona relativo al nodo `gradonum`, no en coordenadas absolutas.
  

  \node[anchor=west,text=milcGradeText,text width=11.4cm,align=left] at ([xshift=1.5cm,yshift=0.5cm]gradonum.east)
    {
     {\sffamily\bfseries\fontsize{9}{11}\selectfont\MakeUppercase{Territorio interior · Socioemocional}}\\[3mm]
     {\sffamily\bfseries\fontsize{21}{25}\selectfont\setlength{\baselineskip}{27pt}El desastre\\[4pt]no es natural}\\[3.5mm]
     {\sffamily\bfseries\fontsize{15}{18}\selectfont Grado 10° · Momento 8}
    };
\end{tikzpicture}

\vspace*{9.4cm}
\vfill

{\sffamily\bfseries\fontsize{8}{10}\selectfont\textcolor{milcGradeDark}{LO QUE TE LLEVAS HOY}}

\begin{tcolorbox}[enhanced,colback=white,colframe=milcNegro,arc=2mm,boxrule=1.6pt,
  left=5mm,right=5mm,top=4mm,bottom=4mm,before skip=3pt,after skip=12pt,
  fontupper=\RaggedRight\fontsize{11}{15.5}\selectfont]
El mapa de vulnerabilidad de tu cuadra: los once ángulos aplicados a un lugar concreto con evidencia observable, y una propuesta de mitigación que no dependa de que alguien de afuera venga a salvarlos.
\end{tcolorbox}

\begin{tcolorbox}[enhanced,colback=milcGris,colframe=milcGris,arc=2mm,boxrule=0pt,
  left=5mm,right=5mm,top=3.5mm,bottom=3.5mm,after skip=12pt]
{\sffamily\bfseries\fontsize{8}{10}\selectfont\textcolor{milcNegro!55}{ESTA GUÍA ES DE}}\\[3mm]
\begin{tabularx}{\linewidth}{X p{3.2cm} p{3.2cm}}
\linead{\linewidth} & \linead{3.0cm} & \linead{3.0cm}\\[-1mm]
{\fontsize{8.5}{10}\selectfont\textcolor{milcNegro!55}{Nombre completo}} &
{\fontsize{8.5}{10}\selectfont\textcolor{milcNegro!55}{Grupo}} &
{\fontsize{8.5}{10}\selectfont\textcolor{milcNegro!55}{Fecha}}\\
\end{tabularx}
\end{tcolorbox}

\vfill

\noindent\textcolor{milcLinea}{\rule{\linewidth}{1pt}}\\[2mm]
{\sffamily\bfseries\fontsize{9.5}{12}\selectfont Autor: Dr. Álvaro Cárdenas Orozco}\\[1mm]
{\sffamily\fontsize{8.5}{11}\selectfont\textcolor{milcNegro!55}{Metodología MILC · Apertura ancestral · Vulnerabilidad global y lectura crítica del riesgo · Triángulo Dussel-Estoicismo-Floridi}}

\newpage

\section*{Apertura: El desastre no es natural: los once ángulos de la vulnerabilidad, aplicados a tu cuadra}

\begin{softbox}{milcGradeDark}{milcGradeSoft}{Saber ancestral}
Cuando el terremoto de 1994 arrasó Tierradentro, el pueblo \textbf{nasa} no lo nombró «desastre natural»: lo leyó como una \textbf{desarmonía} del territorio. Los médicos tradicionales hablaron de un llamado de atención de la Madre Tierra por cómo se estaban relacionando con ella, y señalaron lo que había pasado antes del sismo: la tala, las quemas, la bonanza de cultivos ilícitos que llegó a sembrarse hasta en sitios sagrados, y lo que ellos mismos llamaron una descomposición social. Eso no es superstición: es \textbf{una teoría del riesgo}. Está diciendo que el daño no lo explica solo el movimiento del suelo, sino el estado en que ese movimiento encontró al territorio y a su gente. Tres días después del sismo, el país creó la Corporación Nasa Kiwe para la reconstrucción, y su primer director vinculó a las comunidades al proceso apoyándose en líderes indígenas que conocían la zona y \textbf{manejaban la lengua y la cosmovisión}.
\end{softbox}

\begin{softbox}{milcTurquesa}{milcGris}{Saber contemporáneo}
Ese primer director era \textbf{Gustavo Wilches-Chaux}, y un año antes había publicado el texto que hoy sostiene esta guía. En \emph{Los desastres no son naturales} (LA RED, 1993) define la \textbf{vulnerabilidad} como «la incapacidad de una comunidad para absorber, mediante el autoajuste, los efectos de un determinado cambio en su medio ambiente, o sea su inflexibilidad o incapacidad para adaptarse a ese cambio». Y llama \textbf{vulnerabilidad global} al resultado de muchos factores que se traban entre sí hasta producir «el "bloqueo" o incapacidad de la comunidad para responder adecuadamente ante la presencia de un riesgo determinado, con el consecuente "desastre"». Para estudiarla propone \textbf{once ángulos}. La idea de fondo es una sola y cambia todo: la amenaza es natural, \textbf{el desastre no}. El sismo lo pone la tierra; el desastre lo veníamos construyendo entre todos.
\end{softbox}

\begin{softbox}{milcMostaza}{milcGris}{Pregunta-puente}
\textit{El mismo hombre que dirigió la reconstrucción del territorio nasa escribió que los desastres no son naturales, y los médicos tradicionales le venían diciendo algo parecido con otras palabras. ¿Y si el 10 de agosto la tierra solo hubiera \textbf{revelado} lo que ya estaba mal repartido desde antes?}
\end{softbox}

\begin{softbox}{milcVerde}{milcGris}{Saber y saber hacer}
Al terminar podrás: (1) \textbf{identificar} en una cuadra concreta señales observables de vulnerabilidad, sin señalar a nadie; (2) \textbf{explicar} la diferencia entre amenaza y vulnerabilidad, y por qué el mismo sismo daña distinto a un lado y otro de la misma calle; (3) \textbf{analizar} un caso real reconociendo varios de los once ángulos actuando a la vez; (4) \textbf{crear} el mapa de vulnerabilidad de tu cuadra con una propuesta que la propia comunidad pueda emprender.
\end{softbox}



\small
\begin{tabularx}{\textwidth}{>{\bfseries}p{3.0cm}Y}
\rowcolor{milcGradePrimary!12}Autor & Dr. Álvaro Cárdenas Orozco\\
DBA & Comprende los fenómenos que afectan a su territorio, regula sus emociones ante ellos y acompaña a otros con empatía (transversal 6°--11°, articulado con la política «Escuela, territorio de vida» del MEN, Resolución 006519 de 2025, y la Ley 1620 de 2013)\\
Referentes & Wilches-Chaux (1993) y LA RED · Corporación Nasa Kiwe (Decreto 1179 de 1994) · NSR-10 · Dussel · Epicteto · Floridi\\
Producto final & El mapa de vulnerabilidad de tu cuadra: los once ángulos aplicados a un lugar concreto con evidencia observable, y una propuesta de mitigación que no dependa de que alguien de afuera venga a salvarlos.\\
\end{tabularx}
\normalsize

\puente{Ya sabes que hay dos maneras de decir lo mismo: la de los médicos tradicionales nasa y la de un libro de 1993. Mira ahora el mapa de la sesión: vas a mirar tu cuadra, a entender por qué el daño no cae parejo y a proponer algo que sí esté al alcance.}

\milcestacion{milcGradeDark}{Ruta MILC: así vamos a aprender}
\begin{tabularx}{\textwidth}{>{\centering\arraybackslash}X>{\centering\arraybackslash}X>{\centering\arraybackslash}X>{\centering\arraybackslash}X}
\phasecard{1}{milcVerde}{milcGris}{Escucha\\[-1mm]\small Observo, escucho y pregunto.} &
\phasecard{2}{milcTurquesa}{milcGris}{Sistematización\\[-1mm]\small Analizo el riesgo.} &
\phasecard{3}{milcMagenta}{milcGris}{Praxis\\[-1mm]\small Construyo y aplico.} &
\phasecard{4}{milcVino}{milcGris}{Evaluación\\[-1mm]\small Reflexión liberadora.}
\end{tabularx}


\puente{Antes de la teoría, la calle. Vas a mirar un tramo pequeño de tu barrio con una pregunta muy concreta: si mañana vuelve a temblar, ¿qué de lo que veo aquí decide quién sale bien librado?}

\milcestacion{milcVerde}{Estación 1 de 3 · Escucha}

\begin{softbox}{milcVerde}{milcGris}{Escena para escuchar}
\textbf{Escoge una cuadra} (la tuya, la del colegio, una que recorras a diario) y camínala con una libreta. No estás buscando culpables ni casas «mal hechas»: estás registrando \textbf{hechos observables} de un lugar. Mira el terreno: ¿es plano, es ladera, hay un cauce cerca, hay lleno o relleno? Mira lo construido: materiales, antigüedad aparente, ampliaciones sobre losas que no parecen previstas, muros agrietados. Mira lo colectivo: por dónde saldría la gente, si hay un espacio abierto cercano, si el paso está obstruido. Y anota una cosa más, que casi nadie anota: \textbf{qué se dice} en esa cuadra sobre el riesgo. \textbf{Regla de trabajo, sin excepción:} no fotografíes personas ni fachadas identificables, no escribas nombres ni direcciones exactas, y describe el problema, nunca al dueño. Una casa vulnerable no es una casa culpable.
\end{softbox}

{\sffamily\bfseries\fontsize{8}{10}\selectfont\textcolor{milcNegro!55}{MARCA CUANDO LO HAYAS HECHO}}
\milccheckrow{Delimitaste un tramo concreto y anotaste al menos ocho observaciones, cada una con su ubicación aproximada.}
\milccheckrow{Distinguiste lo que tiene que ver con el terreno de lo que tiene que ver con lo construido.}
\milccheckrow{Registraste al menos una frase que se dice en esa cuadra sobre el riesgo, sin atribuirla a una persona identificable.}

\begin{infoband}{milcVerde}


\textbf{En tu cuaderno (Actividad 1 · IDENTIFICA).} Título: «Actividad 1. Mi cuadra con libreta». \textbf{Formato:} croquis a mano del tramo, más una lista numerada de observaciones ubicadas en el croquis. \textbf{Extensión:} mínimo ocho observaciones y una frase escuchada. \textbf{Sabes que terminaste cuando} tu croquis no permite identificar la casa de nadie y aun así alguien que no conoce el barrio entendería dónde están los problemas.
\end{infoband}

\puente{Ya tienes tus observaciones. Ahora las herramientas para ordenarlas, porque lo que viste no es un montón de problemas sueltos: son ángulos de una misma cosa.}

\milcestacion{milcTurquesa}{Fase 2 · Sistematización: los once ángulos de la vulnerabilidad}

\begin{softbox}{milcTurquesa}{milcGris}{Cuatro claves para leer por qué el daño no cayó parejo}
La frase que da nombre a esta guía suena a provocación y es una definición técnica. La \textbf{amenaza} es el fenómeno: que el suelo se mueva, que el río crezca. Eso es natural y no se negocia. El \textbf{desastre} es lo que ese fenómeno encuentra: casas, materiales, dinero, organización, creencias. Y eso se construye durante años, con decisiones humanas. Wilches-Chaux propuso mirarlo desde once ángulos; aquí los agrupamos en tres familias para leerlos más fácil, advirtiendo que el agrupamiento es nuestro y no del autor.
\end{softbox}

\milcpaso{1}{milcTurquesa}{\textbf{Amenaza no es vulnerabilidad, y la diferencia lo explica todo.} Vulnerabilidad es «la incapacidad de una comunidad para absorber, mediante el autoajuste, los efectos de un determinado cambio en su medio ambiente». Fíjate en lo que no dice: no habla del sismo, habla de la comunidad. Por eso dos municipios con la misma sacudida terminan distinto, y por eso preguntarse «¿por qué a ellos sí?» tiene respuesta, y no es la mala suerte.}
\milcpaso{2}{milcTurquesa}{\textbf{Lo que se ve en el terreno.} Los ángulos \emph{natural}, \emph{físico} y \emph{ecológico}. En El Cairo, el 100 por ciento de las edificaciones sufrió daño y el municipio quedó, en palabras de la gobernadora, «en un 80 por ciento destrozado». Su alcalde explicó por qué: buena parte de las viviendas supera el siglo y está levantada en \textbf{bahareque}. Aquí entra un dato que conviene tener claro: la norma sismorresistente colombiana rige la construcción \textbf{nueva y formal}. Una casa centenaria o autoconstruida nunca se diseñó bajo ninguna norma.}
\milcpaso{3}{milcTurquesa}{\textbf{Quién puede responder.} Los ángulos \emph{económico}, \emph{social}, \emph{político} e \emph{institucional}. Reforzar una casa cuesta plata que muchas familias no tienen, y ese es, según el propio autor, el factor más pesado de todos. Pero no es solo plata: cuenta si la comunidad está organizada, si tiene capacidad real de decidir sobre lo que le afecta, y si las entidades responden o se estorban entre ellas. Una cuadra organizada es literalmente menos vulnerable que una cuadra sola.}
\milcpaso{4}{milcTurquesa}{\textbf{Lo que se sabe y se cree.} Los ángulos \emph{técnico}, \emph{ideológico}, \emph{cultural} y \emph{educativo}. Aquí vive la frase más costosa del país: «aquí nunca ha pasado nada». También vive lo contrario, la creencia de que no vale la pena prepararse porque lo que pase ya está decidido. Y vive esta clase: si la escuela no enseña a leer el propio territorio, está produciendo vulnerabilidad educativa mientras dicta contenidos.}

\begin{drawbox}{milcTurquesa}{Cómo se usan los once ángulos sin volverlos una lista muerta}
\begin{pasoslista}
\item \textbf{Toma un caso real, no un ejemplo inventado.} Una casa, una escuela, una cuadra que exista. Los once ángulos sirven para mirar algo concreto; sueltos, son una tabla que nadie usa.
\item \textbf{Busca al menos tres ángulos actuando a la vez.} Ese es el ejercicio de verdad: una vivienda en bahareque (física) que su dueña no puede reforzar (económica) en un municipio cuyas entidades no llegaron (institucional). Ninguno de los tres explica el daño por sí solo.
\item \textbf{Pregunta cuál de los tres se puede mover.} No todos ceden igual. La vulnerabilidad física de una casa centenaria es cara de corregir; la social y la educativa, no tanto. Ahí está la propuesta que sí se puede hacer.
\end{pasoslista}
\end{drawbox}

\begin{softbox}{milcMostaza}{milcGris}{Errores comunes y su corrección}
\textbf{Reducir todo a la vulnerabilidad física:} es la más visible y la más cara de arreglar, y quedarse ahí lleva a concluir que no hay nada que hacer. \textbf{Confundir vulnerabilidad con culpa:} nadie escoge nacer en la casa que le tocó, y un mapa que señala familias es un mapa que hace daño. \textbf{Creer que la norma sismorresistente cubre todo:} rige lo nuevo y formal, no lo que ya está en pie. \textbf{Decir «fue la naturaleza» y cerrar el tema:} la naturaleza puso la sacudida; el resto lo pusimos nosotros. \textbf{Usar los once ángulos como lista de chequeo:} si no encuentras varios actuando juntos, no estás analizando.
\end{softbox}

\begin{infoband}{milcTurquesa}
\guiaDiagrama{assets/10-8/once-angulos.tikz}{Míralos buscando cuáles coinciden sobre un mismo caso. Un problema explicado por un solo ángulo casi siempre está mal analizado; los que importan aparecen de a tres.}

\textbf{En tu cuaderno (Actividad 2 · ANALIZA).} Título: «Actividad 2. Tres ángulos sobre un caso». \textbf{Formato:} elige \textbf{un} caso de tus observaciones y analízalo con al menos tres de los once ángulos, nombrándolos con el número de la figura. \textbf{Extensión:} un párrafo por ángulo, con la evidencia que lo sostiene. \textbf{Sabes que terminaste cuando} puedes decir cuál de los tres ángulos es el más fácil de mover, y por qué.
\end{infoband}

\puente{Ya sabes leer la vulnerabilidad. Es hora de convertir esa lectura en un mapa que sirva, y en una propuesta que no dependa de que alguien de afuera venga a salvarlos.}

\milcestacion{milcMagenta}{Fase 3 · Praxis: el mapa de la cuadra}

\begin{softbox}{milcMagenta}{milcGris}{Construyendo el mapa de vulnerabilidad con evidencia observable}
Un mapa de vulnerabilidad no es un mapa de miedos ni una denuncia contra los vecinos: es una herramienta para decidir por dónde empezar. Y termina en una propuesta, porque un diagnóstico que no propone nada deja a la gente peor que antes: sabiendo que está en riesgo y sin nada que hacer al respecto.
\end{softbox}

\milcpaso{1}{milcMagenta}{\textbf{Delimitar:} escojan un tramo pequeño y concreto, del tamaño que puedan recorrer y sostener con evidencia. Una cuadra bien mirada vale más que un barrio entero imaginado.}
\milcpaso{2}{milcMagenta}{\textbf{Observar:} vuelvan con la libreta y completen. Cada punto del mapa debe apoyarse en algo que se vio, no en algo que se supone. Si no lo vieron, va marcado como «por confirmar».}
\milcpaso{3}{milcMagenta}{\textbf{Clasificar:} asignen a cada punto los ángulos que lo explican, con la numeración de la figura. Un punto puede llevar varios; los que llevan uno solo casi siempre están mal analizados.}
\milcpaso{4}{milcMagenta}{\textbf{Proponer:} una sola propuesta, realista, que dependa de la comunidad o de la escuela y no de un presupuesto que nadie controla. Con responsable, plazo y forma de saber si funcionó.}

\begin{drawbox}{milcMagenta}{Lo que el mapa debe cumplir antes de mostrarse}
\begin{checklista}
\item \textbf{Ningún punto señala a una familia.} Se describe la condición, nunca al dueño; sin fotos de personas ni fachadas identificables, sin nombres, sin direcciones exactas.
\item \textbf{Cada punto tiene su evidencia}, y lo que no se pudo verificar aparece dicho como «por confirmar».
\item \textbf{Al menos tres puntos} están explicados con dos o más ángulos, con su numeración.
\item \textbf{Una propuesta viable}, con responsable, plazo y una forma concreta de saber si sirvió.
\item \textbf{Está dicho lo que el mapa no puede hacer:} no reemplaza un estudio estructural. Si ven riesgo de colapso, eso se informa a la autoridad municipal, no se queda en un cuaderno.
\end{checklista}
\end{drawbox}

\begin{softbox}{milcMagenta}{milcGris}{Plantilla de trabajo}
\textbf{Plantilla del mapa.} Una hoja apaisada. \emph{Centro:} el croquis del tramo, con puntos numerados. \emph{Margen derecho:} por cada punto, tres renglones — qué se observó, qué ángulos lo explican (con su número) y qué tan fácil sería moverlo (alto, medio, bajo). \emph{Pie de página:} la propuesta única, con responsable y plazo; la frase «este mapa no reemplaza un estudio estructural»; y quién lo hizo y cuándo. \textbf{Sin colores de alarma}: nada de rojo intenso sobre casas de gente real. Un mapa que asusta no se usa, se guarda.
\end{softbox}

\begin{infoband}{milcMagenta}


\textbf{En tu cuaderno (Actividad 3 · CREA).} Título: «Actividad 3. El mapa de mi cuadra». \textbf{Formato:} la hoja completa según la plantilla. \textbf{Extensión:} mínimo seis puntos numerados, tres de ellos con dos o más ángulos, y la propuesta. \textbf{Sabes que terminaste cuando} podrías mostrárselo a alguien que viva en esa cuadra sin que se sienta señalado, y aun así entienda por dónde empezar.
\end{infoband}

\puente{Recorriste y clasificaste. Ahora deja tu evidencia: el mapa, que es tu cuadra vista con once ojos en vez de con uno.}

\milcestacion{milcMostaza}{Producto de la sesión}
\textbf{Mapa de vulnerabilidad de la cuadra, con sus ángulos y una propuesta viable}.
\begin{softbox}{milcMostaza}{milcGris}{Criterios de calidad}
(1) El mapa describe condiciones observables y en ningún punto identifica a una familia o una vivienda concreta. (2) Al menos tres puntos están explicados con dos o más de los once ángulos, nombrados con su numeración. (3) Cada punto tiene evidencia de lo observado, y lo no verificado aparece marcado como tal. (4) La propuesta es una sola, depende de la comunidad o de la escuela, y tiene responsable, plazo y forma de verificar si funcionó. (5) El mapa dice explícitamente lo que no puede hacer, y remite a la autoridad municipal ante sospecha de riesgo estructural.
\end{softbox}

\puente{Terminaste el producto. Detente a mirar qué cambió en ti (no la nota, sino tu manera de mirar la calle por la que pasas todos los días).}

\milcestacion{milcVino}{Evaluación liberadora · cinco dimensiones}

\begin{softbox}{milcVino}{milcGris}{Sentido de la evaluación}
Evaluar no es solo revisar una nota. Es reconocer qué aprendí, qué cambió en mí, qué emociones manejé, cómo me relaciono con la ciudad y la comunidad, y qué le dejaré a quien viene detrás.
\end{softbox}

\textbf{Mi reflexión liberadora en cinco dimensiones.} Completa cada frase iniciada:

\small
\begin{tabularx}{\textwidth}{>{\bfseries\RaggedRight\arraybackslash}p{3.4cm}Y>{\RaggedRight\arraybackslash}p{4.6cm}}
\rowcolor{milcVino}\color{white}Dimensión & \color{white}\bfseries Mi respuesta (frame starter) & \color{white}\bfseries Algo que descubrí\\
1. Personal & Lo que más cambió en mí fue \linead{6.5cm} & \linead{4.2cm}\\
2. Emocional & La emoción más fuerte fue \linead{2.5cm} y la regulé \linead{3cm} & \linead{4.2cm}\\
3. Ciudadana & En democracia esto importa porque \linead{6cm} & \linead{4.2cm}\\
4. Local (Cartago) & En mi barrio sería útil para \linead{6.5cm} & \linead{4.2cm}\\
5. Intergeneracional & A mi abuela/abuelo le diría \linead{6.5cm} & \linead{4.2cm}\\
\end{tabularx}
\normalsize

\begin{infoband}{milcVino}
\textbf{Para la próxima sustentación.} Vuelve a esta tabla en seis meses. Verás que las respuestas cambian — no porque lo aprendido cambie, sino porque tú cambias. La evaluación liberadora es una conversación contigo a través del tiempo.
\end{infoband}


\puente{Cerramos pensando en grande: tres voces para entender por qué lo injusto se disfraza de natural, qué creencias nos dejan quietos y cómo se reparte la información sobre el riesgo.}

\milcestacion{milcGradeDark}{Triángulo de pensamiento · cierre filosófico}

\begin{softbox}{milcVino}{milcGris}{Dussel · lente del nosotros}
\textit{Es una lógica asesina e inmoral, que pasa por ser la naturaleza de las cosas.}

{\footnotesize\itshape\color{milcNegro!70}--- formulación de esta guía, al modo de Enrique Dussel. No es una cita textual.}

Dussel habla de la lógica que ordena el mundo y se disfraza de orden natural, para que nadie pregunte quién lo dispuso así. Que las casas más frágiles sean las de las familias más pobres no es un hecho de la naturaleza: es el resultado de decisiones tomadas durante décadas. Llamarlo «desastre natural» es dejar esas decisiones fuera de la conversación.

\textbf{¿Qué cosa de mi barrio he aceptado siempre como «así es esto», sin preguntarme nunca quién lo decidió?}
\end{softbox}
\linead{\linewidth}\\[1mm]
\linead{\linewidth}

\begin{softbox}{milcMostaza}{milcGris}{Estoicismo · Epicteto · cuidado interior}
\textit{Los hombres se ven perturbados no por las cosas, sino por las opiniones sobre las cosas.}

{\footnotesize\itshape\color{milcNegro!70}--- formulación de esta guía, al modo de Epicteto. No es una cita textual.}

Vale en los dos sentidos, y ese es el filo. Una opinión puede llenarte de miedo sin motivo, y otra puede dejarte quieto cuando había motivo: «aquí nunca ha pasado nada» es una opinión, no un dato, y es de las que más caro cuestan. Wilches-Chaux le da nombre técnico a eso: vulnerabilidad ideológica.

\textbf{¿Qué creo yo sobre por qué pasan estas cosas, y qué hago (o dejo de hacer) por creerlo?}
\end{softbox}
\linead{\linewidth}\\[1mm]
\linead{\linewidth}

\begin{softbox}{milcTurquesa}{milcGris}{Floridi · lente de la infoesfera}
\textit{Somos organismos informacionales (inforgs) conectados entre sí e inmersos en un entorno informacional: la infoesfera.}

{\footnotesize\itshape\color{milcNegro!70}--- formulación de esta guía, al modo de Luciano Floridi. No es una cita textual.}

La información sobre el riesgo también está mal repartida. Quien sabe que su barrio está sobre un lleno, que su casa no cumple norma o que existe un plan municipal, puede decidir; quien no lo sabe, no. Por eso hacer un mapa y devolverlo a la cuadra no es un trabajo escolar: es corregir, en un tramo pequeño, una desigualdad informacional.

\textbf{¿Quién en mi barrio tiene la información sobre el riesgo, y por qué no la tenemos todos?}
\end{softbox}
\linead{\linewidth}\\[1mm]
\linead{\linewidth}

\begin{infoband}{milcNegro}
\textbf{Nota para el docente.} Las tres formulaciones son del autor de la guía, no citas textuales de los pensadores: recogen su lente para pensar el tema, no sus palabras. Si vas a citarlos en clase, remite a la obra original.
\end{infoband}

\milcestacion{milcVino}{Mi autoevaluación · marca una casilla por fila}

{\small
\setlength{\extrarowheight}{2.2pt}
\begin{tabularx}{\textwidth}{>{\RaggedRight\arraybackslash\bfseries}p{3.0cm}YYY}
\rowcolor{milcVino}\color{white}\bfseries Criterio & \color{white}\bfseries Lo logré & \color{white}\bfseries En proceso & \color{white}\bfseries Necesito apoyo\\[.6mm]
Comprensión del tema & \milcopcion{Explico las ideas con mis palabras.} & \milcopcion{Reconozco las ideas, falta precisión.} & \milcopcion{Necesito repasar lo básico.}\\[.8mm]
\rowcolor{milcGris}Lectura del riesgo & \milcopcion{Distingo amenaza de vulnerabilidad y reconozco varios ángulos actuando a la vez en un caso real.} & \milcopcion{Reconozco la vulnerabilidad física, pero aún reduzco todo a la calidad de la construcción.} & \milcopcion{Necesito releer la fase 2.}\\[.8mm]
Aplicación y producto & \milcopcion{Mi producto cumple los criterios.} & \milcopcion{Está armado, con ajustes pendientes.} & \milcopcion{Aún está en construcción.}\\[.8mm]
\rowcolor{milcGris}Reflexión 5 dimensiones & \milcopcion{Respondí cada una con honestidad.} & \milcopcion{Respondí algunas con detalle.} & \milcopcion{Mis respuestas son muy generales.}\\[.8mm]
Triángulo de pensamiento & \milcopcion{Conecté las 3 voces con mi trabajo.} & \milcopcion{Conecté al menos una.} & \milcopcion{No las relaciono con mi proceso.}\\[.8mm]
\end{tabularx}}

\vspace{3mm}
\textbf{Hoy entendí que:}\\[1mm]
\linead{\linewidth}\\[1mm]
\linead{\linewidth}

\textbf{Mi compromiso para la próxima sesión es:}\\[1mm]
\linead{\linewidth}\\[1mm]
\linead{\linewidth}

\begin{softbox}{milcMostaza}{milcGris}{Cita de despedida · Marco Aurelio}
\textit{``El obstáculo en el camino se vuelve el camino. Lo que estorba al camino se vuelve camino.''}\\[1mm]
Cada error de hoy, bien observado, se convierte en pieza del camino que pisarás mañana.
\end{softbox}

\small
\begin{tabularx}{\textwidth}{>{\bfseries}p{3.6cm}Y}
\rowcolor{milcGradeDark!12}Nombre completo: & \linead{10cm}\\
Fecha: & \linead{4cm} \quad Curso/grupo: \linead{4cm}\\
Mi firma: & \linead{9cm}\\
\end{tabularx}
\normalsize

\begin{infoband}{milcGradeDark}
\textbf{Para la próxima sesión.} Dos cosas. Primero, \textbf{devuelve el mapa}: muéstraselo a alguien que viva en esa cuadra, escucha lo que te corrija y anótalo. Un mapa hecho sobre la gente y no con la gente casi siempre está mal, y quien vive ahí sabe cosas que tú no viste en una tarde. Segundo, averigua una sola cosa concreta: si tu municipio tiene plan de gestión del riesgo y dónde se consulta. No para leerlo entero, sino para saber que existe y quién responde por él. Trae una frase: qué te corrigieron del mapa, y qué encontraste.
\end{infoband}

\end{document}
